%==============================================================
%  lecons/algebre/exponentielle_matrice.tex
%  Exponentielle de matrice
%==============================================================

\lecontitle{Exponentielle de matrice}{Algèbre linéaire — Séries de matrices et systèmes différentiels}

%=============================================================
%  § 1 — DÉFINITION ET CONVERGENCE
%=============================================================
\secnum{navyblue}{1}{Définition et convergence}

\begin{tcolorbox}[thmbox]
  \textbf{Définition.}\enspace%
  \index{exponentielle de matrice}
  Pour $A \in \mathcal{M}_n(\mathbb{K})$, on définit
  \[
    e^A = \exp(A) = \sum_{k=0}^{+\infty} \frac{A^k}{k!}
    = I_n + A + \frac{A^2}{2!} + \frac{A^3}{3!} + \cdots
  \]

  \textbf{Théorème.}\enspace%
  \textit{Cette série converge absolument pour toute matrice $A$,
  uniformément sur tout borné de $\mathcal{M}_n(\mathbb{K})$.}
\end{tcolorbox}

\rubriq{Preuve de la convergence}

\begin{mdframed}[style=proofbar]
  On munit $\mathcal{M}_n(\mathbb{K})$ d'une norme sous-multiplicative $\|\cdot\|$
  vérifiant $\|AB\|\leq\|A\|\,\|B\|$ (par exemple la norme subordonnée à la norme
  euclidienne). Alors $\|A^k\|\leq\|A\|^k$, d'où
  \[
    \sum_{k=0}^{N}\left\|\frac{A^k}{k!}\right\| \leq \sum_{k=0}^{N}\frac{\|A\|^k}{k!}
    \xrightarrow[N\to\infty]{} e^{\|A\|} < +\infty.
  \]
  La série est absolument convergente ; comme $\mathcal{M}_n(\mathbb{K})$ est
  un espace de Banach (espace vectoriel normé de dimension finie, donc complet),
  toute série absolument convergente y est convergente, et donc cette série
  converge.

  Enfin, sur toute boule $\{A : \|A\|\leq R\}$, chaque terme est majoré par
  $\|A^k/k!\|\leq R^k/k!$ indépendamment de $A$, le majorant $\sum R^k/k! = e^R$
  étant fini : par le critère de Weierstrass, la convergence est normale, donc
  uniforme, sur tout borné. $\blacksquare$
\end{mdframed}

\rubriq{Premières valeurs}
\[
  e^{0_n} = I_n, \qquad
  e^{\mathrm{diag}(\lambda_1,\ldots,\lambda_n)} = \mathrm{diag}(e^{\lambda_1},\ldots,e^{\lambda_n}).
\]

\decsep{}

%=============================================================
%  § 2 — PROPRIÉTÉS FONDAMENTALES
%=============================================================
\secnum{navyblue}{2}{Propriétés fondamentales}

\begin{tcolorbox}[thmbox]
  Soient $A, B \in \mathcal{M}_n(\mathbb{K})$.
  \begin{enumerate}
    \item \textbf{Commutativité.}\enspace{}Si $AB = BA$, alors $e^{A+B} = e^A e^B$.
    \item \textbf{Inversibilité.}\enspace{}$e^A$ est toujours inversible et
          ${(e^A)}^{-1} = e^{-A}$.
    \item \textbf{Conjugaison.}\enspace{}Si $P$ est inversible,
          $e^{PAP^{-1}} = P\,e^A\,P^{-1}$.
    \item \textbf{Transposée.}\enspace{}$e^{A^T} = {(e^A)}^T$.
    \item \textbf{Déterminant (formule de Jacobi).}\enspace%
    \index{formule de Jacobi}
          $\det(e^A) = e^{\mathrm{tr}(A)}$.
    \item \textbf{Dérivation.}\enspace{}$\dfrac{d}{dt}e^{tA} = A\,e^{tA} = e^{tA}A$.
  \end{enumerate}
\end{tcolorbox}

\rubriq{Preuve de 1 — formule du produit}

\begin{mdframed}[style=proofbar]
  Si $AB=BA$, le binôme de Newton s'applique : ${(A+B)}^k = \sum_{j=0}^k \binom{k}{j}A^j B^{k-j}$.
  Les séries définissant $e^A$ et $e^B$ étant absolument convergentes, leur
  produit de Cauchy converge vers le produit $e^A e^B$ ; on reconnaît alors
  précisément le développement de $e^{A+B}$ :
  \[
    e^{A+B} = \sum_{k=0}^\infty \frac{{(A+B)}^k}{k!}
    = \sum_{k=0}^\infty \sum_{j=0}^k \frac{A^j B^{k-j}}{j!(k-j)!}
    = \left(\sum_{j=0}^\infty\frac{A^j}{j!}\right)\!\left(\sum_{l=0}^\infty\frac{B^l}{l!}\right)
    = e^A e^B.
  \]
  $\blacksquare$
\end{mdframed}

\rubriq{Preuve de 5 — formule de Jacobi}

\begin{mdframed}[style=proofbar]
  \textit{Cas diagonalisable.} Si $A = P\,\mathrm{diag}(\lambda_i)\,P^{-1}$, alors
  $e^A = P\,\mathrm{diag}(e^{\lambda_i})\,P^{-1}$ et
  $\det(e^A) = \prod e^{\lambda_i} = e^{\sum\lambda_i} = e^{\mathrm{tr}(A)}$.

  \textit{Cas général.} La fonction $t\mapsto f(t) = \det(e^{tA})$ vérifie
  $f(0)=1$. Posons $M(t) = e^{tA}$ : on a $M'(t) = A\,e^{tA}$ et $M(t)$ est
  inversible avec $M(t)^{-1} = e^{-tA}$. La formule de Jacobi infinitésimale
  donne, pour $M$ inversible,
  \[
    \frac{d}{dt}\det M(t) = \det M(t)\cdot\mathrm{tr}\bigl(M(t)^{-1}M'(t)\bigr).
  \]
  Ici $M(t)^{-1}M'(t) = e^{-tA}A\,e^{tA} = A$, d'où
  $f'(t) = \det(e^{tA})\,\mathrm{tr}(A) = \mathrm{tr}(A)\,f(t)$.
  Donc $f(t) = e^{t\,\mathrm{tr}(A)}$, et en $t=1$ : $\det(e^A) = e^{\mathrm{tr}(A)}$.
  $\blacksquare$
\end{mdframed}

\decsep{}

%=============================================================
%  § 3 — CALCUL EFFECTIF
%=============================================================
\secnum{navyblue}{3}{Calcul effectif}

\rubriq{Matrice diagonalisable}
Si $A = P\,\mathrm{diag}(\lambda_1,\ldots,\lambda_n)\,P^{-1}$, alors
\[
  e^A = P\,\mathrm{diag}(e^{\lambda_1},\ldots,e^{\lambda_n})\,P^{-1}.
\]

\rubriq{Matrice nilpotente}
Si $A^m = 0$, la série est finie :
\[
  e^A = I + A + \frac{A^2}{2!} + \cdots + \frac{A^{m-1}}{(m-1)!}.
\]

\rubriq{Décomposition de Dunford}
\index{décomposition de Dunford}
Toute matrice $A$ à polynôme caractéristique scindé (en particulier toute
matrice à coefficients dans $\mathbb{C}$) s'écrit de manière unique
$A = D + N$ avec $D$ diagonalisable, $N$ nilpotente et $DN = ND$. Comme $D$ et
$N$ commutent :
\[
  e^A = e^D \cdot e^N,
\]
où $e^D$ se calcule via la diagonalisation et $e^N$ par la série finie.

\begin{mdframed}[style=exbar]
  \textbf{Exemple — matrice $2\times 2$.}\enspace%
  Calculer $e^A$ pour $A = \begin{pmatrix}1&1\\0&1\end{pmatrix} = I + N$,
  $N = \begin{pmatrix}0&1\\0&0\end{pmatrix}$ (nilpotente, $N^2=0$).

  $I$ et $N$ commutent, donc $e^A = e^I\,e^N = e\bigl(I+N\bigr)
  = e\begin{pmatrix}1&1\\0&1\end{pmatrix}$.

  \medskip
  \textbf{Exemple — rotation.}\enspace%
  $A = \theta J$ avec $J = \begin{pmatrix}0&-1\\1&0\end{pmatrix}$.
  Valeurs propres complexes $\pm i\theta$. Comme $J^2 = -I$, on a
  $A^{2k} = (-1)^k\theta^{2k}I$ et $A^{2k+1} = (-1)^k\theta^{2k+1}J$ ; en
  regroupant les termes pairs et impairs de la série, on reconnaît les
  développements de $\cos\theta$ et $\sin\theta$ :
  \[
    e^A = \cos\theta\,I + \sin\theta\,J
        = \begin{pmatrix}\cos\theta & -\sin\theta\\\sin\theta & \cos\theta\end{pmatrix}.
  \]
\end{mdframed}

\decsep{}

%=============================================================
%  § 4 — APPLICATION : SYSTÈMES DIFFÉRENTIELS LINÉAIRES
%=============================================================
\secnum{navyblue}{4}{Application aux systèmes différentiels}

\begin{tcolorbox}[thmbox]
  \textbf{Théorème.}\enspace%
  \index{système différentiel linéaire}
  \textit{Soit $A \in \mathcal{M}_n(\mathbb{R})$. L'unique solution du
  problème de Cauchy}
  \[
    \begin{cases} X'(t) = A\,X(t) \\ X(0) = X_0 \in \mathbb{R}^n \end{cases}
  \]
  \textit{est $X(t) = e^{tA}\,X_0$.}

  \medskip
  \textit{Plus généralement, la solution de $X' = AX + B(t)$,
  $X(0)=X_0$ est donnée par la formule de variation des constantes :}
  \[
    X(t) = e^{tA}\,X_0 + \int_0^t e^{(t-s)A}\,B(s)\,ds.
  \]
\end{tcolorbox}

\begin{mdframed}[style=proofbar]
  \textit{Existence.} $\frac{d}{dt}(e^{tA}X_0) = Ae^{tA}X_0$ (propriété 6) et
  $e^{0\cdot A}X_0 = X_0$.

  \textit{Unicité.} Si $Y$ est une autre solution, posons $Z = e^{-tA}Y$.
  Alors $Z' = -Ae^{-tA}Y + e^{-tA}Y' = e^{-tA}(-AY+AY) = 0$, donc $Z$ est
  constante : $Z(t)=Z(0)=X_0$, soit $Y(t)=e^{tA}X_0$. \hfill$\blacksquare$
\end{mdframed}

\rubriq{Interprétation spectrale}
Si $A$ est diagonalisable avec valeurs propres $\lambda_1,\ldots,\lambda_n$
et vecteurs propres $v_1,\ldots,v_n$, la solution générale est
\[
  X(t) = \sum_{i=1}^n c_i\,e^{\lambda_i t}\,v_i.
\]
Les parties réelles $\mathrm{Re}(\lambda_i)$ gouvernent la stabilité. Dans ce
cas diagonalisable, $X(t)\to 0$ pour \emph{toute} condition initiale $X_0$
si et seulement si tous les $\mathrm{Re}(\lambda_i) < 0$.
(Le critère reste valable dans le cas général en raisonnant sur la réduite de
Jordan, les blocs faisant apparaître des facteurs polynomiaux qui ne modifient
pas la conclusion.)

\decsep{}

%=============================================================
%  § 5 — CONTRE-EXEMPLES, PIÈGES, EXERCICES
%=============================================================
\secnum{exgreen}{5}{Pièges et exercices}

\rubriq{Pièges}

\begin{mdframed}[style=cexbar]
  \textbf{$e^{A+B} \neq e^A e^B$ sans commutativité.}\enspace%
  Pour $A=\begin{pmatrix}0&1\\0&0\end{pmatrix}$ et $B=\begin{pmatrix}0&0\\1&0\end{pmatrix}$,
  $AB\neq BA$ et $e^{A+B}\neq e^Ae^B$.
  C'est la formule de Baker--Campbell--Hausdorff\index{Baker--Campbell--Hausdorff}
  qui corrige cela : $e^Ae^B = e^{A+B+\frac{1}{2}[A,B]+\cdots}$.

  \smallskip\noindent
  \textbf{$e^A = I$ n'implique pas $A = 0$.}\enspace%
  Pour $A = 2\pi i \begin{pmatrix}1&0\\0&-1\end{pmatrix}$,
  $e^A = I_2$ mais $A\neq 0$.

  \smallskip\noindent
  \textbf{$\log(e^A e^B) \neq A + B$ en général.}\enspace%
  Le logarithme matriciel (inverse local de l'exponentielle) n'est pas
  un homomorphisme, contrairement au cas scalaire.
\end{mdframed}

\vspace{6pt}
\exolabel

\begin{mdframed}[style=exobar]
  \begin{enumerate}
    \item Calculer $e^{tA}$ pour $A = \begin{pmatrix}0&-1\\1&0\end{pmatrix}$
          en utilisant $A^2 = -I$. Reconnaître le résultat.

    \item Calculer $e^A$ pour $A = \begin{pmatrix}1&1&0\\0&1&1\\0&0&1\end{pmatrix}$
          (matrice de Jordan d'ordre 3 pour la valeur propre 1).

    \item Montrer que $\det(e^A) > 0$ pour toute matrice réelle $A$.
          En déduire que l'exponentielle ne peut pas être surjective sur
          $\mathrm{GL}_n(\mathbb{R})$.

    \item Résoudre le système $X' = AX$ avec $X(0)=\begin{pmatrix}1\\0\end{pmatrix}$
          et $A = \begin{pmatrix}-1&2\\0&-3\end{pmatrix}$.
          Étudier la stabilité ($X(t)\to 0$ ?)

    \item (Formule de Liouville) Soit $X' = AX$ avec $W(t) = \det(X_1(t),\ldots,X_n(t))$
          le wronskien de $n$ solutions. Montrer que $W(t) = W(0)\,e^{\int_0^t\mathrm{tr}(A)\,ds}$.
  \end{enumerate}
\end{mdframed}

\leconfooter{E.\ Laguerre (1867), J.\ Sylvester — G.\ Peano (1888) — omniprésente en théorie des systèmes}
